\chapter{演示系统与软件质量验证}

\section{系统目标与边界}

演示系统的目的，是让课程验收者看到训练结果如何进入单图推理，不是提供临床服务。输入为常见图片文件，输出为两类概率、预测倾向和警示。系统不接入医院信息系统，不处理DICOM元数据，不保存个人身份信息，也不提供治疗建议。

\section{系统结构}

前端负责文件选择、预览、请求和结果展示；Flask后端负责模型加载、图像解析、预处理和推理。启动命令显式指定checkpoint与device。后端创建predictor时读取checkpoint的model name、classes和settings，避免网页硬编码与训练配置漂移。

请求流程为：用户选择文件；浏览器以multipart/form-data上传；后端检查文件字段；PIL解码并转换RGB；执行Resize、ToTensor和Normalize；模型前向输出softmax；返回JSON。页面把PNEUMONIA概率与阈值比较，但同时展示NORMAL概率，避免只显示单个结论。

\section{模型与权重}

公开主ViT权重大小约1.03GB，SHA256以\texttt{78604698f504d8d9}开头，完整值见复现附录。大文件不直接放入普通Git历史，而通过release附件提供。README说明下载位置和哈希检查。

新版严格实验生成多个checkpoint，不全部作为网页默认模型发布。最终演示模型应在报告冻结后选择，并把对应数据协议、seed、阈值和权重哈希写入release manifest。若仍使用旧ViT网页，仅能展示旧官方基线，页面不得暗示它等于新版最佳模型。

\section{异常处理}

接口对缺失文件返回400状态和结构化错误；无法解码图片时不进入模型；device=auto优先CUDA、不可用则CPU。用户上传的任意图片都可能被模型强制分类，因此页面必须提示“仅适用于胸片课程演示”。系统不应根据扩展名直接信任内容，也不应把异常堆栈返回浏览器。

\section{自动测试}

原项目有阈值、校准、数据工具和网页接口测试。审计后新增患者ID解析与跨split检测测试，并增加温度搜索范围测试。当前测试共17项（以最终pytest输出为准）。网页测试使用fake predictor，不需要加载1GB权重，验证上传成功、JSON字段和缺失文件错误。

单元测试通过不代表端到端模型完全复现。最终还需运行真实checkpoint CPU smoke test、数据审计、评价复算和PDF编译。软件验收按分层证据完成，而不是用一个“测试通过”概括所有质量。

\section{复现入口}

README提供环境安装、数据目录、权重下载、网页启动和主评价命令。严格研究流程还需要：生成患者级Kermany划分、运行审计、按seed训练、评价Kermany/RSNA、聚合bootstrap、校准与阈值。完整命令放在附录，结果文件由artifact index定位。

\section{软件边界}

网页使用默认阈值时可能在RSNA等外部来源上产生大量误报。即使切换简单混合模型，训练数据覆盖仍有限。系统没有输入域检测、质量控制、病人年龄核验或不确定性拒识。出于这些缺失，系统只能作为“算法流程可运行”的课程证据，不能作为医疗产品原型宣传。
